% Source-derived chapter 09: Consequences for arithmetic representations
% Original source: canonical-representations-surface-groups
\chapter{Consequences for arithmetic representations}
\label{canonical-representations-surface-groups-chapter-09}
\source{canonical-representations-surface-groups}

\begin{remark}[Source section]
\label{canonical-representations-surface-groups-chapter-09-source}
\source{canonical-representations-surface-groups}
\source{canonical-representations-surface-groups}
This chapter reproduces the corresponding section of the retrieved
LaTeX source, including its definitions, statements, examples, and
proofs. It has not yet been translated into Lean.
\notready
\end{remark}

% The source section title was: Consequences for arithmetic representations
\label{section:arithmetic-applications}
\source{canonical-representations-surface-groups}
The main arithmetic consequence of \autoref{theorem:finite-image} is \autoref{theorem:arithmetic-consequence} below, which verifies a prediction of the Fontaine-Mazur conjecture, as we now explain. 
Throughout this section, we will no longer be working over $\mathbb C$.

\subsection{Application to relative Fontaine-Mazur}
\label{subsection:application-fontaine-mazur}
\source{canonical-representations-surface-groups}
\begin{definition}
\source{canonical-representations-surface-groups}
\notready
	\label{definition:arithmetic}
\source{canonical-representations-surface-groups}
	Let $X$ be a variety over a finitely-generated field $K$ with algebraic closure $\overline{K}$, and 
$$\rho: \pi_1^{\text{\'et}}(X_{\overline{K}})\to
\on{GL}_r(\overline{\mathbb{Q}}_\ell)$$ a
continuous representation. We
say that $\rho$ is \emph{arithmetic} if its conjugacy class has finite orbit
under the action of $\on{Gal}(\overline{K}/K)$ on the set of isomorphism classes of
$\pi_1^{\text{\'et}}(X_{\overline{K}})$-representations. Here the action is
induced by the natural outer action of $\on{Gal}(\overline{K}/K)$ on
$\pi_1^{\text{\'et}}(X_{\overline{K}})$. 
\end{definition}

The relative form of the Fontaine-Mazur conjecture predicts
that all semisimple arithmetic representations are of geometric origin
\cite[Conjecture 1 bis]{petrov2020geometrically}. Our main arithmetic result, a
straightforward application of \autoref{theorem:finite-image}, is that this is true for representations of low rank on the \emph{generic} $n$-pointed curve of genus $g$. 
In fact, we show such representations are not only of geometric origin:
they even have finite image.
\begin{theorem}
\source{canonical-representations-surface-groups}
\notready\label{theorem:arithmetic-consequence}
\source{canonical-representations-surface-groups}
Let $K$ be a finitely-generated field of characteristic zero and $(C, x_1,
\cdots, x_n)$ a smooth $n$-pointed geometrically connected curve of genus $g$ 
over $K$,
such that the corresponding map $\on{Spec}(K)\to \mathscr{M}_{g,n, \mathbb{Q}}$ factors through the generic point. 
If $r<\sqrt{g+1}$, any continuous arithmetic representation 
\begin{align*}
\rho: \pi_1^{\text{\'et}}(C_{\overline{K}}\setminus\{x_1, \cdots, x_n\})\to \on{GL}_r(\overline{\mathbb{Q}}_\ell) 
\end{align*}
has finite image.
\end{theorem}
We prove this at the end of \autoref{subsection:application-fontaine-mazur}
after spelling out
some consequences.
\begin{remark}
\source{canonical-representations-surface-groups}
\notready\label{remark:esnault-kerz}
\source{canonical-representations-surface-groups}
	Let $K, (C, x_1, \cdots, x_n)$ be as in the theorem, and choose any
	embedding $K\hookrightarrow \mathbb{C}$. Let $C^{\on{an}}\setminus\{x_1,
	\cdots, x_n\}$ be the associated Riemann surface. By \cite[Lemma 7.5.1]{LL:geometric-local-systems}, 
	\autoref{theorem:arithmetic-consequence} implies that arithmetic
	representations are not Zariski-dense in the
	$\overline{\mathbb{Q}}_\ell$-points of the character variety
	parametrizing semisimple representations of
	$\pi_1(C^{\on{an}}\setminus\{x_1, \cdots, x_n\})$ of rank $r$, for
	$1<r<\sqrt{g+1}$. This answers negatively a question of Esnault and Kerz
	\cite[Question 9.1(1)]{esnault2020arithmetic}.
\end{remark}


Recall that a local system $\mathbb V$ on a smooth complex variety $X$ is of geometric origin if there is
some Zariski open $U \subset X$ and a smooth proper morphism $f: Y \to U$ so that
$\mathbb V|_U$ is a subquotient of $R^i f_* \mathbb C$ for some $i$.
\begin{corollary}
\source{canonical-representations-surface-groups}
\notready
	\label{corollary:fontaine-mazur}
\source{canonical-representations-surface-groups}
	Let $K$ and $(C, x_1, \cdots, x_n)$ be as in
	\autoref{theorem:arithmetic-consequence}, 
	and let
	$\rho$ be a continuous representation 
\begin{align*}
\rho: \pi_1^{\text{\'et}}(C_{\overline{K}}\setminus\{x_1, \cdots, x_n\})\to \on{GL}_r(\overline{\mathbb{Q}}_\ell) 
\end{align*}
with $\dim \rho < \sqrt{g+1}$. The following are equivalent.
	\begin{enumerate}
		\item $\rho$ has finite image.
		\item $\rho$ is arithmetic.
		\item For any embedding $K\hookrightarrow \mathbb{C}$ and any isomorphism $\mathbb{C}\simeq \overline{\mathbb{Q}}_\ell$, the local system corresponding to $\rho_\mathbb{C}|_{C^{\on{an}}\setminus\{x_1,
	\cdots, x_n\}}$ on $C^{\on{an}}\setminus\{x_1,
	\cdots, x_n\}$ is of geometric
			origin. 
		\item For any embedding $K\hookrightarrow \mathbb{C}$ and any isomorphism $\mathbb{C}\simeq \overline{\mathbb{Q}}_\ell$, the local system corresponding to $\rho_\mathbb{C}|_{C^{\on{an}}\setminus\{x_1,
	\cdots, x_n\}}$ on $C^{\on{an}}\setminus\{x_1,
	\cdots, x_n\}$ underlies an
			integral PVHS.
	\end{enumerate}
	Moreover, $(1), (3)$, and $(4)$ are equivalent whenever $\dim \rho <
	2\sqrt{g+1}$.
\end{corollary}
\begin{proof}
The equivalence of $(1)$ and $(2)$ when $\dim \rho <\sqrt{g+1}$ holds by
	\autoref{theorem:arithmetic-consequence}, and $(1)$
	certainly implies $(3)$ and $(4)$ for any value of $\dim \rho$.
	Next, $(3)$ implies $(4)$ for any value of $\dim \rho$ since any local system of geometric origin
	underlies an integral PVHS. This is well known, but see, for example, the proof of
	\cite[Corollary 1.2.7]{LL:geometric-local-systems}.

	To conclude, it remains to show that $(4)$ implies $(1)$ when $\dim \rho
	< 2\sqrt{g+1}$.
	This follows from \cite[Corollary 1.2.7]{LL:geometric-local-systems},
	by choosing an embedding $K\hookrightarrow \mathbb{C}$ corresponding to an analytically very general complex point of $\mathscr{M}_{g,n}$.
	By \cite[Corollary 1.2.7]{LL:geometric-local-systems},
	 a local system on an analytically very general curve underlying an integral PVHS has finite monodromy whenever its dimension is
	$< 2\sqrt{g+1}$.
	Therefore, $\rho$ also has finite monodromy.
\end{proof}
\begin{remark}
\source{canonical-representations-surface-groups}
\notready\label{remark:fontaine-mazur}
\source{canonical-representations-surface-groups}
	The relative Fontaine-Mazur conjecture predicts that all semisimple arithmetic  local systems
	are of geometric origin. Hence on a generic curve, all semisimple arithmetic local systems $\rho$ with $\dim \rho< 2 \sqrt{g+1}$ should  have finite monodromy by
	\autoref{corollary:fontaine-mazur}.
	In particular, this suggests one should be able to improve the bound of $\sqrt{g+1}$ in
	\autoref{theorem:arithmetic-consequence} to $2 \sqrt{g+1}$ if one
	restricts to semisimple representations.

Note that arithmeticity of $\rho$ on $C$ over $K$ is an a stronger condition than having finite orbit under
the action of the mapping class group after base change to $\mathbb C$, and so
this does not mean the bound of \autoref{theorem:finite-image-semisimple} can be improved
to $2 \sqrt{g+1}$. Indeed, it cannot be improved, even in the semisimple case, when $g = 1$, see
\autoref{remark:sharp-g-1}.
%	In general, one expects that local systems of geometric origin are the same as arithmetic local systems on one hand, and local systems underlying a polarizable integral variation of Hodge structure on the other. Combining \cite[Corollary 1.2.7]{LL:geometric-local-systems} with \autoref{theorem:arithmetic-consequence} shows that these three notions indeed agree for local systems of rank $<\sqrt{g+1}$ on the generic $n$-pointed curve of genus $g$.
\end{remark}


%\subsubsection{}
%\label{subsubsection:arithmetic-proof}
\begin{proof}[Proof of \autoref{theorem:arithmetic-consequence}]
	Let $C^\circ_{\overline K} := 
	C_{\overline{K}}\setminus\{x_1, \cdots, x_n\}$. For $R$ a ring, let $\mathscr
	M_{g,n,R}$ denote the moduli stack of genus $g$, $n$-pointed curves over
	$\on{Spec} R$
	and let $\pi^\circ_R: \mathscr C_R^\circ \to \mathscr M_{g,n,R}$ denote the universal $n$-punctured genus
	$g$ curve over $R$.
	We use $C^\circ_R$ to denote a fiber of $\pi^\circ_R$.

	We claim (suppressing basepoints) that there is a map of exact sequences of \'etale fundamental groups
	\begin{equation}
	\label{equation:map-of-birman-exact-sequences}
\source{canonical-representations-surface-groups}
	\begin{tikzcd}
	0 \ar {r} & \pi^\text{\'et}_1(C^\circ_{\overline K}) \ar {r} \ar {d}{\simeq} &
	\pi^\text{\'et}_1(C^\circ_K) \ar {r} \ar {d} & \pi^\text{\'et}_1(\on{Spec} K) \ar {r}
	\ar {d} & 0 \\
	0 \ar {r} & \pi^\text{\'et}_1(C^\circ_{\overline {K(M_{g,n,\mathbb Q})}}) \ar {r} &
	\pi^\text{\'et}_1(\mathscr C^\circ_{\mathbb Q}) \ar {r} & \pi^\text{\'et}_1(\mathscr
	M_{g,n,\mathbb Q}) \ar {r} & 0.
	\end{tikzcd}\end{equation}
	Indeed, the first sequence is the standard homotopy exact sequence, see
	\cite[\href{https://stacks.math.columbia.edu/tag/0BTX}{Tag
	0BTX}]{stacks-project}.
	The second one is less standard but, using 
	\cite[\href{https://stacks.math.columbia.edu/tag/0BTX}{Tag
	0BTX}]{stacks-project} again, its exactness
	can be reduced to verifying exactness of
	\begin{equation}
	\label{equation:q-curve-fibration}
\source{canonical-representations-surface-groups}
	\begin{tikzcd}
	0 \ar {r} &  \pi^\text{\'et}_1(C^\circ_{\overline {K(M_{g,n,\mathbb Q})}}) \ar {r} &
	\pi^\text{\'et}_1(\mathscr C^\circ_{\overline {\mathbb Q}}) \ar {r} & \pi^\text{\'et}_1(\mathscr
	M_{g,n,\overline{ \mathbb Q}}) \ar {r} & 0.
	\end{tikzcd}\end{equation}
	By comparison of fundamental groups of algebraically closed fields of
	characteristic $0$, it is enough to verify exactness of the sequence
	with all fields in subscripts replaced by the complex numbers.
	The version for topological fundamental groups is given in
	\autoref{lemma:mgn-fundamental-group}. Since the \'etale fundamental
	group is the profinite completion of the topological fundamental group
	\cite[Expos\'e XII, Corollaire 5.2]{sga1},
	and profinite completions of exact sequences where the left term has
	trivial center remain exact \cite[Proposition
	3]{anderson:exactness-properties}, we obtain exactness of
	\eqref{equation:q-curve-fibration}.

We next use \eqref{equation:map-of-birman-exact-sequences} to deduce that 
$\rho$ has finite orbit under the natural action of $\pi_1(\mathscr
M_{g,n,\mathbb \overline{\mathbb Q}})$, to be constructed below.
	The two exact sequences in
	\eqref{equation:map-of-birman-exact-sequences} induce compatible outer actions of
	$\pi_1(\on{Spec} K)$ and $\pi_1(\mathscr
	M_{g,n,\mathbb Q})$ on 
	$\pi_1(C^\circ_{\overline {K(M_{g,n,\mathbb Q})}})$.
	Hence, the outer action of
	$\pi_1(\on{Spec} K)$ factors through that of $\pi_1(\mathscr
	M_{g,n,\mathbb Q})$.
	Because $\on{Spec} K \to \mathscr
	M_{g,n,\mathbb Q}$ is dominant, the induced map
$\pi_1(\on{Spec} K) \to \pi_1(\mathscr M_{g,n,\mathbb Q})$ has image of finite
index.
In particular, the orbit of the conjugacy class of $\rho$ under
$\pi_1^{\text{\'et}}(\mathscr{M}_{g,n, \mathbb{Q}}, \overline{\eta})$,
and hence under $\pi_1^{\text{\'et}}(\mathscr{M}_{g,n,
\overline{\mathbb{Q}}}, \overline{\eta})$, is finite. 

	Choose an embedding $\overline{K}\hookrightarrow\mathbb{C}$, and let
	$(\mathscr{C}^\circ)^{\text{an}}$ be the corresponding punctured Riemann
	surface, so that,
	by residual finiteness of surface groups,
	there is a natural inclusion
	$\pi_1((\mathscr{C}^\circ)^{\text{an}})\hookrightarrow
	\pi_1^{\text{\'et}}(\mathscr{C}^\circ_{\overline{K}})$. As the conjugacy
	class of $\rho$ has finite orbit under
	$\pi_1^{\text{\'et}}(\mathscr{M}_{g,n, \overline{\mathbb{Q}}},
	\overline{\eta})
	\simeq \pi_1^{\text{\'et}}(\mathscr{M}_{g,n, \mathbb{C}},
\overline{\eta}_{\mathbb C})$, the standard comparison theorems between \'etale and
	topological $\pi_1$
	\cite[Expos\'e XII, Corollaire 5.2]{sga1}
	imply that
	$\rho|_{\pi_1((\mathscr{C}^\circ)^{\text{an}})}$ has finite orbit under
	the topological fundamental group
	$\pi_1(\mathscr{M}_{g,n, \mathbb{C}}^{\text{an}},
	\overline{\eta}_{\mathbb{C}})$. Hence
	$\rho|_{\pi_1((\mathscr{C}^\circ)^{\text{an}})}$ is MCG-finite, and hence has
	finite image by \autoref{theorem:finite-image},
	which may be applied after choosing an isomorphism between $\overline{\mathbb{Q}}_\ell$ and $\mathbb{C}$. 
	But
	$\pi_1( (\mathscr{C}^\circ)^{\text{an}})$ is dense in
	$\pi_1^{\text{\'et}}(\mathscr{C}^\circ_{\overline{K}})$, so $\rho$ also
	has finite image.
%	We first reduce to verifying this in the case that $K$ is the function
%	field of $\mathscr M_{g,n}$.
%	Let $L$ be the function field of $\mathscr{M}_{g,n}/\mathbb{Q}$, so that
%	$K/L$ is a finitely generated extension. Let $\overline{K}$ be an
%	algebraic closure of $K$ and let $\overline{L}$ be the algebraic closure
%	of $L$ in $\overline{K}$. Let $\mathscr{C}^\circ/L$ be the generic curve
%	with the $n$ tautological points removed, so that (by the definition of $\mathscr{M}_{g,n}$) there is a canonical
%	isomorphism $C\setminus\{x_1, \cdots, x_n\}\simeq\mathscr{C}^\circ_K$,
%	inducing an isomorphism on geometric fundamental groups. The induced map
%	$\on{Gal}(\overline{K}/K)\to \on{Gal}(\overline{L}/L)$ has image of
%	finite index, and the outer action of $\on{Gal}(\overline{K}/K)$ on
%	$\pi_1^{\text{\'et}}(C\setminus\{x_1, \cdots,
%	x_n\}_{\overline{K}})\simeq
%	\pi_1^{\text{\'et}}(\mathscr{C}^\circ_{\overline{L}})$ factors through
%	this map. So viewing $\rho$ as a representation of
%	$\pi_1^{\text{\'et}}(\mathscr{C}^\circ_{\overline{L}})$, it is  
%	arithmetic. Thus, it suffices to consider the case of the generic curve over the function field of $\mathscr{M}_{g,n}$.
%	
%	Now $\mathscr{C}^\circ$ has good reduction over all of
%	$\mathscr{M}_{g,n}$; thus the action of $\on{Gal}(\overline{L}/L)$ on
%	$\pi_1^{\text{\'et}}(\mathscr{C}^\circ_{\overline{L}})$ factors through
%	the natural surjection $\on{Gal}(\overline{L}/L)\to
%	\pi_1^{\text{\'et}}(\mathscr{M}_{g,n, \mathbb{Q}}, \overline{\eta})$,
%	where $\overline{\eta}$ is the geometric generic point of
%	$\mathscr{M}_{g,n}$ corresponding to our choice of $\overline{L}$. In
%	particular, the orbit of the conjugacy class of $\rho$ under
%	$\pi_1^{\text{\'et}}(\mathscr{M}_{g,n, \mathbb{Q}}, \overline{\eta})$,
%	and hence under $\pi_1^{\text{\'et}}(\mathscr{M}_{g,n,
%	\overline{\mathbb{Q}}}, \overline{\eta})$, is finite. (Here
%	$\overline{\mathbb{Q}}$ is the algebraic closure of $\mathbb{Q}$ in $\overline{L}$.)
%	
%	Now choose an embedding $\overline{L}\hookrightarrow\mathbb{C}$, and let
%	$(\mathscr{C}^\circ)^{\text{an}}$ be the corresponding punctured Riemann
%	surface, so that,
%	by residual finiteness of surface groups,
%	there is a natural inclusion
%	$\pi_1((\mathscr{C}^\circ)^{\text{an}})\hookrightarrow
%	\pi_1^{\text{\'et}}(\mathscr{C}^\circ_{\overline{L}})$. As the conjugacy
%	class of $\rho$ has finite orbit under
%	$\pi_1^{\text{\'et}}(\mathscr{M}_{g,n, \overline{\mathbb{Q}}},
%	\overline{\eta})
%	\simeq \pi_1^{\text{\'et}}(\mathscr{M}_{g,n, \mathbb{C}},
%\overline{\eta}_{\mathbb C})$, the standard comparison theorems between \'etale and
%	topological $\pi_1$
%	\cite[Expos\'e XII, Corollaire 5.2]{sga1}
%	imply that
%	$\rho|_{\pi_1((\mathscr{C}^\circ)^{\text{an}})}$ has finite orbit under
%	the topological fundamental group
%	$\pi_1(\mathscr{M}_{g,n, \mathbb{C}}^{\text{an}},
%	\overline{\eta}_{\mathbb{C}})$. Hence
%	$\rho|_{\pi_1((\mathscr{C}^\circ)^{\text{an}})}$ is MCG-finite, and hence has
%	finite image by \autoref{theorem:finite-image}
%	which may be applied after choosing an isomorphism between $\overline{\mathbb{Q}}_\ell$ and $\mathbb{C}$. 
%	But
%	$\pi_1( (\mathscr{C}^\circ)^{\text{an}})$ is dense in
%	$\pi_1^{\text{\'et}}(\mathscr{C}^\circ_{\overline{L}})$, so $\rho$ also
%	has finite image.
\end{proof}

\subsection{Application to lifting residual representations}
Let $K, (C, x_1, \cdots, x_n)$ be as in
\autoref{theorem:arithmetic-consequence}. We next explain why there are residual
representations of $\pi_1^{\text{\'et}}(C_{\overline{K}}\setminus\{x_1, \cdots,
x_n\})$ which do not lift to arithmetic representations in characteristic zero.
In particular, in \autoref{example:non-liftable-representations}
we construct residual representations of geometric
fundamental groups which are not ``of geometric origin," as we now define. To our knowledge these are the first such examples.
\begin{definition}
\source{canonical-representations-surface-groups}
\notready
Let $K$ be an algebraically closed field, and let $X/K$ be a variety. Let $\mathbb{F}$ be a finite field of characteristic different from that of $K$. Let $L$ be the fraction field of the Witt vectors $W(\mathbb{F})$. We say that a continuous representation $$\rho: \pi_1^{\text{\'et}}(X_{\overline{K}})\to \on{GL}_r(\mathbb{F})$$ is \emph{of geometric origin} if there exists a continuous representation $$\xi: \pi_1^{\text{\'et}}(X_{\overline{K}})\to \on{GL}_r(\overline{L})$$ such that:
\begin{enumerate}
\item There exists a $\xi$-stable
	$\mathscr{O}_{\overline{L}}$-lattice $W\subset \overline{L}^r$ with
	$W\otimes \overline{\mathbb{F}}\simeq \rho \otimes
	\overline{\mathbb{F}}$.
\item There exists a dense open subscheme $U\subset X$ and a smooth proper morphism $\pi: Y\to U$, such that $\xi|_{\pi_1^{\text{\'et}}(U_{\overline{K}})}$ arises as a subquotient of the monodromy representation of $R^i\pi_*\overline{L}$ for some $i\geq 0.$
\end{enumerate}
\end{definition}
In other words, we say a residual representation is of geometric origin if it
arises as the reduction of a characteristic zero representation of geometric
origin.
\begin{corollary}
\source{canonical-representations-surface-groups}
\notready\label{corollary:lifting-and-geometric-origin}
\source{canonical-representations-surface-groups}
	Let $K, (C, x_1, \cdots, x_n)$ be as in \autoref{theorem:arithmetic-consequence}. Then for $1<r<\sqrt{g+1}$ and  $p\gg_r 0$, no surjective representation 
	$$\rho: \pi_1^{\text{\'et}}(C_{\overline{K}}\setminus\{x_1, \cdots, x_n\})\to \on{GL}_r(\mathbb{F}_p)$$ admits an arithmetic lift to characteristic zero. In particular no such representation is of geometric origin.
\end{corollary}
\begin{proof}
	By definition, a mod $p$ representation of geometric origin lifts to a
representation of geometric origin over $\overline{\mathbb{Q}}_p$. As representations of geometric origin are arithmetic, the second statement follows from the first.
	
	Thus it suffices to show that for $p\gg 0$ and $\rho$ as in the theorem, $\rho$ does not admit an arithmetic lift. Suppose to the contrary that it did admit an arithmetic lift $\xi$; by \autoref{theorem:arithmetic-consequence}, $\xi$ would have finite image. Thus it suffices to show that for $Q$ a finite, totally ramified extension of $\mathbb{Q}_p$, there do not exist any finite subgroups of $\on{GL}_r(\mathscr{O}_Q)$ surjecting onto $\on{GL}_r(\mathbb{F}_p)$ if $p\gg_r 0$. 
	
	This follows from Jordan's theorem on finite subgroups of $\on{GL}_r(Q)$,
	where $Q$ is any field of characteristic zero
	(\cite[p.~91]{jordan1878memoire} or \cite[Theorem
	36.13]{curtis1966representation}). Recall that Jordan's theorem says
	that there exists some constant $n(r)$ such that if $G\subset \on{GL}_r(Q)$ is a finite subgroup, $G$ contains an abelian normal subgroup of index dividing $n(r)$. In particular, for any $g_1, g_2\in G$, we have $$[g_1^{n(r)}, g_2^{n(r)}]=I_r,$$ where $I_r$ is the $r\times r$  identity matrix.
	
	Now define $\overline{g_1}, \overline{g_2}\in \on{GL}_r(\mathbb{F}_p)$ by
	$$\overline{g_1}=\begin{pmatrix} 1 & 1 \\ 0 & 1 \end{pmatrix} \oplus
	I_{r-2}$$ $$\overline{g_2}=\begin{pmatrix} 1 & 0 \\ 1 & 1 \end{pmatrix}
	\oplus I_{r-2}.$$ Direct computation shows that $[\overline{g_1}^{n(r)},
	\overline{g_2}^{n(r)}]\neq I_r$ as long as $p$ does not divide $n(r)$.
	But if $\on{GL}_r(\mathbb{F}_p)$ was the image of a finite subgroup of
	$\on{GL}_r(\mathscr{O}_Q)$, this equality would hold by Jordan's theorem, giving the claim. Indeed, we have shown the result for all $p$ not dividing the constant $n(r)$ from Jordan's theorem.
\end{proof}
Using \autoref{corollary:lifting-and-geometric-origin}, we now construct
examples of residual representations with no arithmetic lifts.
\begin{example}
\source{canonical-representations-surface-groups}
\notready
	\label{example:non-liftable-representations}
\source{canonical-representations-surface-groups}
	Fix non-negative integers $g,n$ so that $n \geq 1$ if $g =1$ and $n\geq
	3$ if $g = 0$.
	Fix $r$ with $1<r<\sqrt{g+1}$. We claim
	that for any $n$-pointed genus $g$ curve $(C, x_1, \cdots, x_n)$, and any $p$,
	there exist surjective representations $$\rho:
	\pi_1^{\text{\'et}}(C_{\overline{K}}\setminus\{x_1, \cdots, x_n\})\to
	\on{GL}_r(\mathbb{F}_p)$$ as in \autoref{corollary:lifting-and-geometric-origin}. 
	In particular these representations do not admit arithmetic lifts to
	characteristic $0$.

	If $n > 0$, maintaining the assumption $n \geq 3$ when $g = 0$, we have
	that
$\pi_1^{\text{\'et}}(C_{\overline{K}}\setminus\{x_1, \cdots, x_n\})$ is
profinite free on at least $2$ generators.
The claim in this case follows from the fact that $\on{GL}_r(\mathbb{F}_p)$ is generated by two elements
\cite{waterhouse:two-generators}. 

We now deal with the case $n =0$ and $g \geq 2$.
Indeed, this follows from the fact that
$\pi_1^{\text{\'et}}(C_{\overline{K}})$ surjects onto a free profinite group on
$g$ generators by writing it as the profinite completion of $\langle a_1, \ldots,
a_g, b_1, \ldots, b_g \rangle/ \prod_{i=1}^g [a_i, b_i]$ and considering the map
to the free profinite group generated by $c_1, \ldots, c_g$ sending $a_i \mapsto
c_i, b_i \mapsto \on{id}$.
\end{example}

We conclude the section with several remarks describing consequences of the
above non-liftable residual representations.
\begin{remark}
\source{canonical-representations-surface-groups}
\notready\label{remark:dejong}
\source{canonical-representations-surface-groups}
\autoref{corollary:lifting-and-geometric-origin} provides a counterexample to a particularly optimistic extension of de Jong's conjecture \cite[Conjecture 2.3, Theorem 3.5]{de2001conjecture}.	De Jong's conjecture, proven by Gaitsgory \cite{gaitsgory2007jong}, implies that an absolutely irreducible residual $\mathbb{F}_q$-representations of geometric fundamental groups of smooth curves over finite fields of characteristic not dividing $q$ always lift to arithmetic representations over a field of characteristic zero. One might naturally ask if the same is true for curves over arbitrary finitely-generated fields; \autoref{corollary:lifting-and-geometric-origin} shows that this is not the case.
\end{remark}
\begin{remark}
\source{canonical-representations-surface-groups}
\notready
	\label{remark:complete-intersection}
\source{canonical-representations-surface-groups}
	The solution to de Jong's conjecture, as described in
	\autoref{remark:dejong}, implies that deformation rings of
absolutely irreducible $\mathbb{F}_p$-representations of arithmetic fundamental
groups of curves over finite fields (of characteristic different from $p$) are
always complete intersections over $\mathbb{Z}_p$. Our results show that the
analogous statement is not true for the arithmetic fundamental group
$\pi_1^{\text{\'et}}(C_{K}\setminus\{x_1, \cdots, x_n\})$ as in
\autoref{theorem:arithmetic-consequence} because they are not flat over $\mathbb
Z_p$ by \autoref{example:non-liftable-representations}. 
\end{remark}
\begin{remark}
\source{canonical-representations-surface-groups}
\notready
	\label{remark:flach}
\source{canonical-representations-surface-groups}
	Flach asks  \cite[p.~7]{cornelissen2005problems} if deformation rings of
	absolutely irreducible residual representations of profinite groups are
	always complete intersections. By now it is well-known that the answer
	to this question is in general ``no" (see e.g.~\cite{eardley2016inverse}
	for a more or less complete answer to this question, and the references
therein). Our result \autoref{corollary:lifting-and-geometric-origin} shows that
this question has a negative answer even for arithmetic fundamental groups of
generic smooth  curves, as explained in \autoref{remark:complete-intersection}.
\end{remark}
